\documentclass[pra,twocolumn,superscriptaddress,longbibliography]{revtex4-2}
\usepackage[a4paper, left=0.8in, right=0.8in, top=1.0in, bottom=1.0in]{geometry}

\usepackage{amssymb}
\usepackage{amsthm}
\usepackage{amsmath}
\usepackage{url}
\usepackage{graphicx}

\usepackage[braket, qm]{qcircuit}

\usepackage{mathtools}
\DeclarePairedDelimiter\ceil{\lceil}{\rceil}
\DeclarePairedDelimiter\floor{\lfloor}{\rfloor}

\newtheorem{definition}{Definition}
\newtheorem{theorem}{Theorem}
\newtheorem{corollary}[theorem]{Corollary}
\newtheorem{lemma}[theorem]{Lemma}
\newtheorem{proposition}[theorem]{Proposition}
\newtheorem{problem}{Problem}
\theoremstyle{remark}
\newtheorem*{remark}{Remark}


\newcommand{\ea}[1]{\begin{align}#1\end{align}}
\newcommand{\eat}[2]{\begin{alignat}{#1} #2\end{alignat}}
\newcommand{\nn}{\nonumber \\}
\newcommand{\la}{\langle}
\newcommand{\ra}{\rangle}
\newcommand{\mcl}{\mathcal}
\newcommand{\mbf}{\mathbf}
\newcommand{\mbb}{\mathbb}
\newcommand{\mrm}{\mathrm}
\newcommand{\bmp}{\begin{pmatrix}}
\newcommand{\emp}{\end{pmatrix}}
\newcommand{\vpu}[1]{^{\vphantom{#1}}}
\DeclareMathOperator{\tr}{Tr}
\DeclareMathOperator{\vect}{vec}
\DeclareMathOperator{\ebr}{ebr}
\DeclareMathOperator{\len}{len}
\DeclareMathOperator{\ch}{\mcl{C}}
\DeclareMathOperator{\rank}{Rank}
\DeclareMathOperator{\GL}{GL}
\DeclareMathOperator{\Gal}{Gal}
\DeclareMathOperator{\rel}{Re}
\DeclareMathOperator{\sgn}{Sign}
\DeclareMathOperator{\argmax}{argmax}
\DeclareMathOperator{\var}{\mathrm{Var}}
\DeclareMathOperator{\lsh}{LSH}
\DeclareMathOperator{\mom}{\mcl{M}}

\newcommand{\kket}[1]{|#1\rangle\!\rangle}
\newcommand{\bbra}[1]{\langle\!\langle#1|}
\newcommand{\bbrakket}[2]{\langle\!\langle#1 \rvert #2\rangle\!\rangle}

\numberwithin{equation}{section}

\begin{document}

\title{Learnable Quantum Projections for state-vector similarity search}
\author{Sergii Strelchuk and Danylo Yakymenko}
\date{}

\begin{abstract}
We propose an algorithm for solving the state-vector similarity search problem: 
given a database of pure quantum states with exact classical descriptions and multiple copies of a pure query state, 
return a state from the database that has approximately maximal fidelity with the query state. 
The algorithm is based on selecting a set of projections represented by shallow quantum circuits with trainable parameters, 
computing locality-sensitive hashes from projections' outputs for each state in the database, 
and learning the parameters that help the hashes better discriminate among the states in the database. 
% In principle, the algorithm could work for both very large databases (billions of states) and very small ones, comprising only a few states, with an arbitrary amount of qubits. 
We present theoretical grounds for this approach along with the results of our simulation experiments using this method. 
The results indicate that it could be a practical solution for use on currently available NISQ devices in some scenarios.

{\it Keywords}: 
vector similarity search, vector database, approximate nearest neighbor, locality-sensitive hashing,
random projection, dimensionality reduction, quantum machine learning, quantum state discrimination.
\end{abstract}

\maketitle

\section{Introduction}
In the era of rising AI technologies, databases of high-dimensional vectors have become ubiquitous 
due to their ability to encode any kind of information, 
whether it is text, images, media, or any structured data. 
Finding vectors similar to a given one in a database of vectors is thus the cornerstone of information retrieval tasks. 
However, this problem becomes increasingly computationally challenging as database sizes and vector dimensionality increase.
During past decades, many efforts have been put onto this problem and numerous methods have been developed. 
One particular technique, known as the locality-sensitive hashing (LSH) \cite{LSH2002}, will be most relevant to our work. 
The main idea of the LSH is to map high-dimensional vectors to strings of bits, 
where each bit is likely to be same for close vectors. 
The Hamming distance between the bit strings then corresponds to the similarity between vectors, 
and one can use the hashes for the search. 
A simple way to construct LSH is based on the use of random projections, 
which define hyperplanes that split the entire vector space into two halves. 
Each hash bit indicates which half a vector falls into. 
We refer to \cite{OptimalLSH2006, AngularLSH2015} as examples of the evolution of this method. 
Given a database of vectors, the selected projections can greatly affect the search performance, 
so one can try to learn how to select them in the optimal way.
This is known as data-dependent hashing, 
the survey \cite{survey2016} gives an overview of such learning techniques.
An interested reader can find a collection of resources devoted to the LSH methods on the webpages \cite{page_Andoni, page_LWJ}. 

Quantum computing, on the other hand, has the potential to make significant contributions to the machine learning field 
due to its ability to represent high-dimensional vectors as quantum states of many-qubit systems and operate on them efficiently. 
The state similarity search problem is also closely related to the quantum state discrimination \cite{QSD2009, QSD2015} 
and quantum tomography in general, 
with applications extending beyond machine learning to areas like quantum cryptography. 
In comparison to the classical case, the quantum case depends on the number of copies of the query state available. 
Thankfully, with the invention of the shadow tomography in \cite{shadow2017}, 
we don't have to fully reconstruct a state to be able to predict many of the values of target observables. 
The shadow tomography was later improved in \cite{shadow2020} (see also the follow-up \cite{shadow2022}) 
where authors proposed the classical shadows algorithm, 
which requires only a logarithmic number of measurements in the number of target observables 
to predict their values up to some error. 
In particular, this allows to predict fidelities between the query state and the states in the database, 
which makes the similarity search possible. 
We will argue, however, that our proposed algorithm could perform better than classical shadows approach for the similarity search problem. 

More specifically, the problem we consider is the following: 

\begin{problem}[Main problem] \label{prob:main}
Let $\mcl{H} = \mbb{C}^d, d = {2^q}$, be a Hilbert space representing a quantum system with $q$ qubits.
Given a database of pure quantum states $\ket{v_1}, \dots, \ket{v_N} \in \mcl{H}$ with exact classical descriptions
and multiple copies of a pure query state $\ket{u}$, 
find a state in the database that has approximately maximal fidelity $|\ip{u}{v_i}|^2$ with the query state. 
Try to minimize the total number of quantum and classical operations used by the algorithm.
\end{problem}

If only a single copy of the query state is given, then this essentially becomes the quantum state discrimination problem 
(even though we do not assume that $\ket{u}$ is exactly one of the states $\ket{v_i}$). 
The ``pretty good measurement'' (PGM) \cite{PGM1994, PGM2007}, defined as a positive operator-valued measure (POVM) with elements 
$\{\rho^{-1/2}\op{v_i}{v_i}\rho^{-1/2}\}$, where $\rho = \frac{1}{N}\sum_i\op{v_i}{v_i}$, 
provides a sufficiently good recipe for the solution. 
However, in general, a POVM implementation requires an exponential number of primitive gates, and it produces only a probabilistic answer.
Thus, PGM may not be optimal in practice, especially when multiple copies of the query state $\ket{u}$ are available. 
In what follows, we do not impose any restrictions on the number of copies of the query state $\ket{u}$, 
although such a restriction is implicit since it affects the overall complexity of the algorithm $-$ which is our main concern.
One can find resources devoted to the multi-copy quantum state discrimination 
(see, e.g., \cite{MQSD2019} and references therein), 
though they mostly focus on schemes that are optimal with regard to the discrimination error, 
independently of the implementation complexity. 

The main idea of our algorithm is based on the following observation:

\begin{theorem}[Main theorem] \label{thm:main}
Let $\ket{v}$ be a state in $\mcl{H}$, 
$P$ be a projection onto some $d_s$ dimensional subspace 
such that $\melem{v}{P}{v} = l$, 
and $\ket{u}$ be a Haar random state with the fixed fidelity $|\ip{v}{u}|^2 = f$. Then 
\begin{align}
& {\mbb{E}} \melem{u}{P}{u} = \frac{d_s}{d} + \bigg(f - \frac{1}{d}\bigg)\bigg(l - \frac{d_s}{d}\bigg)\frac{d}{d-1}, \\
& {\var} \melem{u}{P}{u} \le  \frac{(1-f)d}{4(d-1)^2}\bigg( f + \frac{d - 2}{d}\bigg) \le \frac{1}{4d}, 
\end{align}
where $\mbb{E}$ and $\var$ denote the expectation and variance, respectively.
\end{theorem}

Note that without the constrain $|\ip{v}{u}|^2 = f$, one can find that ${\mbb{E}} \melem{u}{P}{u} = \frac{d_s}{d}$ and 
${\mbb{E}} |\ip{v}{u}|^2 = \frac{1}{d}$. The theorem implies that if $l > \frac{d_s}{d}$ and $f > \frac{1}{d}$ then 
the expectation ${\mbb{E}} \melem{u}{P}{u}$ is greater than $\frac{d_s}{d}$ with exponentially small variance 
(which tends to $0$ with the increase of $f$). 
Thus, the condition $\melem{u}{P}{u} > \frac{d_s}{d}$ becomes a good indicator of the closeness of the states $\ket{v}$ and $\ket{u}$,
and the higher the fidelity $|\ip{v}{u}|^2$ and the length $\melem{v}{P}{v}$ are, the stronger this indicator is. 
Similarly, if $\melem{v}{P}{v} < \frac{d_s}{d}$ with $f > \frac{1}{d}$, then ${\mbb{E}} \melem{u}{P}{u} < \frac{d_s}{d}$. 
In Section \ref{sec:pre} we give an overview of the tools used in the proof of this theorem, with a full proof found in the Appendix. % \ref{sec:proof}. 

The theorem motivates to make the following definition: 
\begin{definition}[Projection based locality-sensitive hash of a state]\label{def:lsh}
Given a projection $P$, the locality-sensitive hash function for a state $\ket{v} \in \mcl{H}$ is defined by
\begin{align}
\lsh_P(\ket{v}) = \begin{cases}
+1, & \text{if } \melem{v}{P}{v} > \frac{d_s}{d} + \tau, \\
-1, & \text{if } \melem{v}{P}{v} < \frac{d_s}{d} - \tau, \\
0, & \text{otherwise},
\end{cases}
\end{align}
where $\tau$ is a small threshold parameter that essentially cuts the cases where $\melem{v}{P}{v}$ is very close to the average.

For a list of projections $P_1,\dots,P_n$, the locality-sensitive hash of $\ket{v}$ is the concatenation of hashes for each projection
\begin{align}
\lsh_{P_1,\dots,P_n}(\ket{v}) = (\lsh_{P_1}(\ket{v}),\dots,\lsh_{P_n}(\ket{v})).
\end{align}
\end{definition}

Since the variance of $\melem{v}{P}{v}$ for a Haar random $\ket{v}$ is also exponentially small, 
most of the time $\lsh_P(\ket{v})$ will be $0$ for a meaningful (e.g., greater than $1/\sqrt{d}$) value of threshold $\tau$. 
However, if the number $N$ of state-vectors $\ket{v_i}$ in a database is large, then, by chance, 
some of them will have $\lsh_P(\ket{v_i})=\pm 1$. And if we take a sufficiently high number of projections $P_1, \dots, P_n$, 
then for each vector in the database, we can expect at least a small number of non-zero LSH values. 
We can vary the threshold $\tau$ to control the number of non-zero LSH values. 
Additionally, we do not have to pick the projections $P_i$ randomly. 
We can select them based on the LSH values they produce for the database vectors, 
or train parametrized projections using an optimization algorithm to maximize the number of non-zero LSH values. 
The training is particularly useful if the database is small. %, consisting of only a few vectors.
Speaking of the choice of rank $d_s$ of projections, we have found that $d_s = d/2 = 2^{q-1}$ suits best from a practical perspective. 
Projections of such a rank can be implemented as unitaries followed by the measurement of a single qubit. 
To be more efficient, we use one unitary circuit to produce $q$ projections of rank $2^{q-1}$ $-$ corresponding to the measurements of each qubit. 

Having the locality-sensitive hash vectors $\lsh_{P_1, \dots, P_n}(\ket{v})$ 
for each state $\ket{v}$ in the database, 
there are multiple ways of how one can use these values for the actual similarity search. 
More details about our algorithm could be found in Section \ref{sec:algo}.  

In general, on NISQ devices, noise adds volatility and could lead to biases in the output expectations of circuit executions. 
While volatility can be addressed by increasing the number of execution shots, an expectation bias is a bigger issue since it cannot be easily predicted. 
This, however, is not a big concern in our algorithm due to its nature $-$ the locality-sensitive hash uses multiple projections, 
and so the possible biases in their executions should not significantly affect the cumulative locality-sensitive property of the hash. 
Moreover, a bias in the execution of a particular circuit, if it exists, should be repeatable. 
This means that if $\lsh_P(\ket{v}) > 0$ for some $\ket{v}$ only due to an execution bias, 
the same hash value should be obtained for states close to $\ket{v}$ since the bias is likely to be the same. 

In Section \ref{sec:res} we present the results of our simulation experiments, 
which indicate that the proposed approach could be a practical solution on NISQ devices in some scenarios. 
In particular, we consider two cases $-$ a database of $100$ thousands $12$-qubit states, 
for which a selection of random projections with no training were used, 
and a database of one thousand $20$-qubit states, for which a set of projections was trained. 
Databases of states on larger amount of qubits have a limited capability in the training of projections since this would not be possible to perform on classical hardware. 
However, we believe that quantum learning of many-qubit projections will become practical on future quantum devices and so this direction deserves further research.

\section{Operator moments}\label{sec:pre}

To prove our main theorem we need basic facts about the computation of operator moments over the Haar measure on the unitary group. 
For a thorough introduction to this subject, we refer readers to \cite{Mele2024}.

Let $\mcl{H} = \mbb{C}^d$ be a Hilbert space of dimension $d$, 
$\mcl{L}(\mcl{H})$ be the space of linear operators on $\mcl{H}$. 
By $\mcl{U}(d)$ we denote the unitary group acting on $\mcl{H}$, 
$\mu$ denotes the Haar measure on $\mcl{U}(d)$. 
The measure $\mu$ also induces a measure on states in $\mcl{H}$ 
by representing them as $U\ket{0}$ for $U \in \mcl{U}(d)$. 

The $k$-th operator moment of an operator $O \in \mcl{L}(\mcl{H}^{\otimes k})$ with respect to the Haar measure is defined by 
\begin{align}
    & \mcl{M}^{(k)}(O) = \mbb{E}_{U \sim \mu} U^{\otimes k}O(U^\dagger)^{\otimes k} \\
    & = \int_{\mcl{U}(d)} U^{\otimes k}O(U^\dagger)^{\otimes k} {\rm d} \mu.
\end{align}

The first operator moment equals to
\begin{align}
\mcl{M}^{(1)}(O) = \tr(O)I/d 
\end{align}
using which one can find that 
\begin{align*}
    & {\mbb{E}}_{\ket{u} \sim \mu} \melem{u}{O}{u}
    = {\mbb{E}}_{U \sim \mu} \tr(UOU^\dagger \op{0}{0})
    = \tr(O)/d, \\
% \end{align*}
% and
% \begin{align*}
    & {\mbb{E}}_{\ket{u},\ket{v} \sim \mu} |\ip{u}{v}|^2 
    = {\mbb{E}}_{\ket{u},\ket{v} \sim \mu} \ip{u}{v}\ip{v}{u} \\
    & = {\mbb{E}}_{U \sim \mu} \tr(U\op{0}{0}U^\dagger\op{0}{0})
    = 1/d. 
\end{align*}

The second operator moment helps in computing the variances of the same quantities. 
Let us denote by $\mbb{I}$ the identity operator on $\mcl{L}(\mcl{H} \otimes \mcl{H})$, 
and by 
\begin{align}
    \mbb{F} = \sum_{i,j} \ket{i}\ket{i}\bra{j}\bra{i}
\end{align}
the flip (swap) operator acting on the same space. 
It is known that 
\begin{align}
    \mcl{M}^{(2)}(O) = c_{\mbb{I}, O} \mbb{I} + c_{\mbb{F}, O}\mbb{F},
\end{align}
where 
\begin{align}
    & c_{\mbb{I}, O} = \frac{\tr(O) - d^{-1}\tr(\mbb{F}O)}{d^2-1}, \\
    & c_{\mbb{F}, O} = \frac{\tr(\mbb{F}O) - d^{-1}\tr(O)}{d^2-1}.
\end{align}

Hence for $O \in \mcl{L}(\mcl{H})$
\begin{align*}
    & {\mbb{E}}_{\ket{u} \sim \mu} (\melem{u}{O}{u})^2 \\
    & = {\mbb{E}}_{U \sim \mu} \tr\big(U^{\otimes 2}O^{\otimes 2}(U^\dagger)^{\otimes 2} (\op{0}{0})^{\otimes 2}\big) \\
    & = \tr( \mcl{M}^{(2)}(O^{\otimes 2}) (\op{0}{0})^{\otimes 2}) \\
    & = c_{\mbb{I}, O^{\otimes 2}} + c_{\mbb{F}, O^{\otimes 2}} = \tr(O^{\otimes 2}\frac{\mbb{I} + \mbb{F}}{d(d+1)}).
\end{align*}

In particular, ${\mbb{E}}_{\ket{u} \sim \mu} |\ip{u}{0}|^4 = \frac{2}{d(d+1)}$. 

TODO: finish section
% Though, it is known that 

\section{The algorithm description}\label{sec:algo}

As was explained in the introduction, our algorithm requires selecting a list of projections $P_1, \dots, P_n$ and 
computing the hashes $\lsh_{P_1, \dots, P_n}(\ket{v_i})$ for each state $\ket{v_i}$ in the database. 
Given multiple copies of a query state $\ket{u}$, we can compute its hash $\lsh_{P_1, \dots, P_n}(\ket{u})$ and use it to 
find the most similar state $\ket{v_i}$ from the database. 
A straightforward way of doing it is by computing the Euclidean distances between the hashes of $\ket{u}$ and $\ket{v_i}$ 
and returning the index $i$ corresponding to the minimal value. A bucketing techniques can be used to make this more efficient. 
Naturally, all vectors $\ket{v_i}$ with the condition $\lsh_{P_1}(\ket{v_i})=1$ fall into the same group, a bucket. If 
$\lsh_{P_1}(\ket{u})=1$ for the query state then it makes sense to search for the most similar state within the same bucket first
(and we can continue the search in other buckets). The use of buckets increases the memory requirements, 
though other forms of bucketing exists that balance the trade-off between time and memory. 
Finally, one could simply train a classical neural network that for a given hash predicts the index of the state in the database that 
has the most similar hash. In essence, one could view the transition from a state to its locality-sensitive hash 
as a form of dimensionality reduction and thus consider any appropriate algorithms on top of it.

The real question in this approach, though, is how do we pick the projections $P_i$ that give us the most representative hashes. 
As was mentioned earlier, picking random projections may not be the best choice. For example, if the database is small, 
let's say it has just two states $\ket{v_1}$, $\ket{v_2}$, then for a random projection $P$ their hashes will be $0$ most of the time 
if the threshold $\tau$ in the $\lsh$ definition is a multiple of the standard deviation (which is exponentially small).
If instead we pick $P$ with learnable parameters then we can train it in way that exactly one of the states $\ket{v_1}$, $\ket{v_2}$ 
has non-zero hash. Then the hash of a query state $\ket{u}$ should tell us what state is closer to $\ket{u}$ with high probability. 

Finding such a projection $P$ may not be possible if the circuits we consider are too shallow.
Clearly, there is a monotonic relationship between the depth of $P$ and the maximum value of $\melem{v}{P}{v}$ one can achieve. 
No depth restriction corresponds to the maximum value of $1$, but this could make our algorithm exponentially complex. 
Moreover, finding depth restricted $P$ is even harder if we want to obtain non-zero hashes for multiple states $\ket{v_i}$ from the database. 
Nevertheless, we think that it is the way to go in selecting $P_i$ because we have to move $\melem{v_j}{P_i}{v_j}$ 
only slightly away from the average and only for some of the states in the database. 
Concerning the above, we define the following main objective in selecting the projections $P_i$:

\begin{definition}[Main objective in selecting $P_i$] \label{def:obj}
Given a database of states $\ket{v_i}$, in the selection of a set of $n$ projections $\{P_i\}$, that has a fixed depth restriction, 
we want to maximize the value 
\begin{align}
\sum_{i,j} |\lsh_{P_i}(\ket{v_j})|.
\end{align}
\end{definition}

Maximizing this objective does not guarantee that two distant vectors will not get the same hashes, but we can use other steps to avoid that possibility. 
We also want states in the database to be covered evenly by non-zero hash values, 
that is, $\sum_{i} |\lsh_{P_i}(\ket{v_j})|$ should be approximately same for each $j$.

In the following subsections, we describe ways of how one could achieve this objective. 
There is one common trick that we use in all of them. 
To not waste the circuit complexity of the projection implementation, 
which is just an application of a unitary $U$ and measuring the first qubit, 
we let each circuit $U$ to produce $q$ projections $\Pi^{(i)}_{U} = U^\dagger \op{0}{0}_i U$, 
each corresponding to measuring the $i$-th qubit after the application of $U$. 

\subsection{Selection of projections without training}

At first, we describe a simple approach where we don't use any training of projections. 
This is useful when the database is quite large having at least $O(d)$ elements. 

The algorithm is as follows:

1. Select the threshold value $\tau$. For example, we might want to have about $10\%$ chance of getting 
$lsh_{P}(\ket{v}) \neq 0$ on average. One can estimate the corresponding $\tau$ by running simulation experiments. 
The diagram \ref{fig:tau} show the estimated values of $\tau$ for the different number of qubits and threshold percent. 

% 1, 2, 4, 8, 16
% q 10 12 14 16 18 20

TODO: rewrite random part

1. Pick a random circuit $U_1$ of the maximal allowed depth and measure all qubits for each $U_1\ket{v_j}$ some number of times (say $1000$) 
to find the estimates $\hat{l}_{1,i,j}$ of the values $l_{1,i,j} = \melem{v_j}{\Pi^{(i)}_{U_1}}{v_j}$. For each qubit $i$, 
sort the values $\hat{l}_{1,i,j}$ and select the maximal and minimal subsets of some cutoff size 
(say $5\%$ of the total number of vectors, each). 
Define $\lsh_{1,i}(\ket{v_j}) = 1$ if $\hat{l}_{1,i,j}$ falls into the maximal subset, 
$\lsh_{1,i}(\ket{v_j}) = -1$ if $\hat{l}_{1,i,j}$ is in the minimal subset, and set it to $0$ otherwise. 
This definition of $\lsh$ is slightly different than we gave before because the threshold $\tau$ is implicit and not constant here, 
but the effect is essentially the same. One could define the average of the implicit values of $\tau$ as some baseline value. 

2. Repeat steps 1 and 2 for some other random unitary $U_2$ and continue to do so until some break condition.  
This condition could be a limit on the number of unitaries or we could look at the total number of non-zero hash values 
and judge on that. 
% 2. Look by how many non-zero hash values each state is covered. 
% Select half of the database states that are covered by least amount of non-zero hash values. 
% Repeat steps 1 and 2 to for this half of database (adjust the cutoff size in step 1 appropriately) for the next random unitary $U_2$ and so on. 

% \quad 

% Since on each iteration after the first one we look only at the half of database states

Note that depending on the problem size, this could be done entirely on a simulator. 

\subsection{The basic training}

\subsection{Adaptive training technique}

\section{Results of experiments}\label{sec:res}

% \appendix

\onecolumngrid

\quad 

\section{Acknowledgments}

\quad

\section{Appendix}

\subsection{Proof of the Main Theorem}\label{sec:proof}

\begin{proof}
Since we are considering the Haar random $\ket{u}$, the distribution of $\melem{u}{P}{u}$ will be the same for pairs 
$(\ket{v'}, P') = (U\ket{v}, UPU\dagger)$ where $U$ is an arbitrary unitary. 
It is not hard to derive (either by using Jordan's lemma about two projectors or simply via step-by-step unitary reductions) 
that any pair $(\ket{v}, P)$ with $\melem{v}{P}{v}=l$, $l \in [0,1]$, 
is unitary equivalent to $(\ket{v_0}, P_0)$ where  
\begin{align*}
& \ket{v_0} = \ket{\bf 0}, ~~P_0 = \bmp P_{2} & 0 \\ 0 & P_{d-2} \emp, \\
& P_{2} = \bmp l & \sqrt{l(1-l)} \\ \sqrt{l(1-l)} & 1-l\emp, ~P_{d-2} = I_{d_s-1} \oplus {\bf 0}_{d-d_s-1}.\\ 
\end{align*}
Without loss of generality, we can prove the theorem only for such pairs $(\ket{v_0}, P_0)$.

Since the fidelity $|\ip{\bf 0}{u}|^2 = f$ is fixed, any such state $\ket{u}$ can be represented as 
$$
\sqrt{f}\ket{\bf 0} + \sqrt{1-f}U\ket{\bf 1},
$$
where %$\ket{v_0^\perp}$ is some fixed unit vector orthogonal to $\ket{\bf 0}$ and 
$U = 1 \oplus U_{d-1}$ is an arbitrary unitary that acts non-trivially only on the subspace orthogonal to $\ket{\bf 0}$.
%with $U_{d-1}\ket{0^q} = \ket{0^q}$. 
% In other words, $U_{d-1}$ acts non-trivially only on the subspace orthogonal to $\ket{0^q}$. 
Thus, we assume that $U_{d-1}$ is a Haar random unitary on this subspace. 
Let's compute the moments over unitaries $U = 1 \oplus U_{d-1}$. 
Assume $O = \bmp O_{00} & O_{01} \\ O_{10} & O_{11} \emp$ is an arbitrary operator on $\mcl{H}$, 
% where $O_{00} \in \mbb{C}^{1\times1}$ and $O_{11} \in \mbb{C}^{(d-1)\times (d-1)}$. 
where $O_{00}$ and $O_{11}$ are matrices of sizes $1\times1$ and $(d-1)\times (d-1)$ respectively. 
Then 
\begin{align*}
    & (1 \oplus U_{d-1}) O (1 \oplus U_{d-1}^\dagger) 
    = \bmp O_{00} & O_{01} U_{d-1}^\dagger \\ U_{d-1} O_{10} & U_{d-1} O_{11} U_{d-1}^\dagger \emp. 
\end{align*}

Since ${\mbb{E}} U_{d-1} = 0$ and ${\mbb{E}}U_{d-1} O_{11} U_{d-1}^\dagger = \tr(O_{11})I\frac{1}{d-1}$, we have that 

\begin{align*}
    & {\mbb{E}}(1 \oplus U_{d-1}) O (1 \oplus U_{d-1}^\dagger) 
    = \bmp O_{00} & 0 \\ 0 & \tr(O_{11})I\frac{1}{d-1} \emp. 
\end{align*}

Now we can compute the expectation of $\melem{u}{P_0}{u}$: 
\begin{align*}
& {\mbb{E}}\melem{u}{P_0}{u} 
= {\mbb{E}}\tr\big(P_0 \op{u}{u} \big)
= \tr\bigg[P_0 \bmp f & 0 \\ 0 & \frac{(1-f)}{d-1}I \emp \bigg] \\
& = lf + (1-l)\frac{(1-f)}{d-1} + (d_s-1)\frac{(1-f)}{d-1} \\
& = \frac{d_s}{d} + \big(f - \frac{1}{d}\big)\big(l - \frac{d_s}{d}\big)\frac{d}{d-1}.
\end{align*}

For the variance we have 
\begin{align*}
& \var\melem{u}{P_0}{u} = \mbb{E}\melem{u}{P_0}{u}^2 - (\mbb{E}\melem{u}{P_0}{u})^2, \\
& \mbb{E}\melem{u}{P_0}{u}^2 
= \tr(P_0^{\otimes2} \cdot \mbb{E}(\op{u}{u})^{\otimes2}).
\end{align*}

%The variance calculation 
This requires us to know how to compute the second reduced moment of an arbitrary operator $O \in \mcl{H}^{\otimes2}$: 
$$
\mom_1^{(2)}(O) = \mbb{E}(1 \oplus U_{d-1})^{\otimes2} O (1 \oplus U_{d-1}^\dagger)^{\otimes2}.
$$  
To simplify analysis we consider the cases where $O$ is a matrix unit. 

1. Let $O = \op{\bf 00}{\bf 00}$. 
Then $\mom_1^{(2)}(O) = \op{\bf 00}{\bf 00}$. 

2. Let $O = \op{\bf 00}{\bf 0j}$ where ${\bf j} \in \{1,\dots,d-1\}$. 
Then
\begin{align*}
& \mom_1^{(2)}(O) = \op{\bf 0}{\bf 0} \otimes \ket{\bf 0}\mbb{E}\bra{\bf j} (1\oplus U_{d-1}^\dagger) = 0. 
\end{align*}
Similarly, we get $0$ for $\op{\bf 00}{\bf j0}$, $\op{\bf j0}{\bf 00}$ and $\op{\bf 0j}{\bf 00}$.

3. Let $O = \op{\bf 00}{\bf ij}$ where ${\bf i,j} \in \{1,\dots,d-1\}$. 
Then 
\begin{align*}
& \mom_1^{(2)}(O) = \ket{\bf 00}\mbb{E}\bra{\bf ij} (1 \oplus U_{d-1}^\dagger)^{\otimes2} = 0. 
\end{align*}
Same for $\op{\bf ij}{\bf 00}$.

4. Let $O = \op{\bf 0i}{\bf 0j}$ where ${\bf i,j} \in \{1,\dots,d-1\}$. 
Then 
\begin{align*}
& \mom_1^{(2)}(O) = \op{\bf 0}{\bf 0} \otimes \mom_1^{(1)}(\op{\bf i}{\bf j}) 
= \op{\bf 0}{\bf 0} \otimes (0\oplus\ip{\bf i}{\bf j}I\frac{1}{d-1}). 
\end{align*}
Similarly for $\op{\bf i0}{\bf j0}$.

5. Let $O = \op{\bf i0}{\bf 0j}$ where ${\bf i,j} \in \{1,\dots,d-1\}$. 
Then 
\begin{align*}
& \mom_1^{(2)}(O) 
= \sum_{{\bf k}=1}^{d-1}\op{\bf k0}{\bf 0k}\frac{\ip{\bf i}{\bf j}}{d-1}. 
\end{align*}
To see this, one can apply the swap operator from the left of $O$, 
which reduces it to the previous case. 
Similarly for $O = \op{\bf 0i}{\bf j0}$. 

6. Let $O = \op{\bf 0i}{\bf jk}$ where ${\bf i,j,k} \in \{1,\dots,d-1\}$. 

Using vectorization we get 
\begin{align*}
& \vect(\mom_1^{(2)}(O)) 
= \mbb{E}(1\oplus\overline{U}_{d-1})^{\otimes2}\otimes I \otimes (1\oplus{U}_{d-1}) \vect(\op{\bf 0i}{\bf jk}) \\
& = \mbb{E}(1\oplus\overline{U}_{d-1})^{\otimes2}\otimes I \otimes (1\oplus{U}_{d-1}) \ket{\bf jk0i} = 0, 
\end{align*}
due to Lemma [?]. The same result we get for 
$\op{\bf i0}{\bf jk}$, $\op{\bf ij}{\bf 0k}$ and $\op{\bf ij}{\bf k0}$. 

7. Let $O = \op{\bf ij}{\bf kl}$ where ${\bf i,j,k,l} \in \{1,\dots,d-1\}$. 

Again, from vectorization and Lemma [?] we get 
\begin{align*}
& \vect(\mom_1^{(2)}(O)) 
% = \mbb{E}(1\oplus\overline{U}_{d-1})^{\otimes2}\otimes (1\oplus{U}_{d-1})^{\otimes2} \ket{\bf klij} \\
= \mbb{E}(1\oplus\overline{U}_{d-1})^{\otimes2}\otimes (1\oplus{U}_{d-1})^{\otimes2} \kket{O} \\
& = \frac{1}{(d-1)^2-1}\left[\kket{\mathbb{I}_{d-1}}\bbrakket{\mathbb{I}_{d-1}}{O}-\frac{1}{d-1}\kket{\mathbb{I}_{d-1}}\bbrakket{\mathbb{F}_{d-1}}{O}
-\frac{1}{d-1}\kket{\mathbb{F}_{d-1}}\bbrakket{\mathbb{I}_{d-1}}{O}+\kket{\mathbb{F}_{d-1}}\bbrakket{\mathbb{F}_{d-1}}{O}\right], 
\end{align*}

Now, $\bbrakket{\mathbb{I}}{O} = \tr(O) = \ip{\bf i}{\bf k}\ip{\bf j}{\bf l}$ and 
$\bbrakket{\mathbb{F}}{O} = \tr(\mbb{F}O) = \ip{\bf i}{\bf j}\ip{\bf k}{\bf l}$. 
Thus 
\begin{align*}
& \mom_1^{(2)}(O)
= \frac{1}{d(d-2)}\left[\mbb{I}_{d-1}(\ip{\bf i}{\bf k}\ip{\bf j}{\bf l} - \frac{\ip{\bf i}{\bf j}\ip{\bf k}{\bf l}}{d-1})
+ \mbb{F}_{d-1}(\ip{\bf i}{\bf j}\ip{\bf k}{\bf l} - \frac{\ip{\bf i}{\bf k}\ip{\bf j}{\bf l}}{d-1})\right].
\end{align*}

In particular, for $O = \op{\bf 11}{\bf 11}$ we get 
\begin{align*}
& \mom_1^{(2)}(\op{\bf 11}{\bf 11})
= \frac{1}{d(d-1)}\left[\mbb{I}_{d-1} + \mbb{F}_{d-1}\right].
\end{align*}

We are now prepared to compute $\mbb{E}(\op{u}{u})^{\otimes2}$:  
\begin{align*}
& \mbb{E}(\op{u}{u})^{\otimes2}
= \mbb{E}(1\oplus{U}_{d-1})^{\otimes2}
\big(f\op{\bf 0}{\bf 0} + (1-f)\op{\bf 1}{\bf 1} + \sqrt{f(1-f)}(\op{\bf 0}{\bf 1} + \op{\bf 1}{\bf 0})\big)^{\otimes2}
(1\oplus{U}_{d-1}^\dagger)^{\otimes2} \\
& = f^2\op{\bf 00}{\bf 00} 
+ \frac{(1-f)^2}{d(d-1)}\left[\mbb{I}_{d-1} + \mbb{F}_{d-1}\right] + \\
& + \frac{f(1-f)}{d-1}\big(\op{\bf 0}{\bf 0} \otimes I_{d-1} + I_{d-1} \otimes \op{\bf 0}{\bf 0} 
+ \sum_{{\bf k}=1}^{d-1}\op{\bf k0}{\bf 0k} + \sum_{{\bf k}=1}^{d-1}\op{\bf 0k}{\bf k0}\big).
\end{align*}

Note that 
\begin{align*}
& P_0^{\otimes2} 
= \bigg(l\op{\bf 0}{\bf 0} + (1-l)\op{\bf 1}{\bf 1} + \sqrt{l(1-l)}(\op{\bf 0}{\bf 1}+\op{\bf 1}{\bf 0}) 
    + \sum_{{\bf j}=2}^{d_s}\op{\bf j}{\bf j} \bigg)^{\otimes2}.
\end{align*}

Then 
\begin{align*}
& \tr(P_0^{\otimes2} \cdot \mbb{E}(\op{u}{u})^{\otimes2}) 
= l^2f^2 + \frac{2(1-l)^2(1-f)^2}{d(d-1)} + \frac{4l(1-l)f(1-f)}{d-1} + \\
& + \frac{2(d_s-1)lf(1-f)}{d-1} + \frac{2(d_s-1)(1-l)(1-f)^2}{d(d-1)} 
+ \frac{(1-f)^2}{d(d-1)}\big( (d_s-1)^2 + (d_s-1)\big) \\
& = l^2f^2 + f(1-f)\bigg(\frac{4l(1-l)}{d-1} + \frac{2(d_s-1)l}{d-1} \bigg) 
+ (1-f)^2\bigg(\frac{2(1-l)^2}{d(d-1)} + \frac{2(d_s-1)(1-l)}{d(d-1)} + \frac{d_s(d_s-1)}{d(d-1)}\bigg) \\
& = l^2f^2 + f(1-f)\frac{l(2d_s+2-4l)}{d-1}
+ (1-f)^2\frac{2(1-l)(d_s-l) + d_s(d_s-1)}{d(d-1)}
\end{align*}

Recall that 
\begin{align*}
& \bigg({\mbb{E}}\melem{u}{P_0}{u}\bigg)^2
= \bigg(lf + (d_s-l)\frac{(1-f)}{d-1} \bigg)^2 
= l^2f^2 + (1-f)^2\frac{(d_s-l)^2}{(d-1)^2} + f(1-f)\frac{l(2d_s-2l)}{d-1}
\end{align*}

Hence 
\begin{align*}
& \var\melem{u}{P_0}{u}
= f(1-f)\frac{2l(1-l)}{d-1}
+ (1-f)^2\bigg(\frac{2(1-l)(d_s-l) + d_s(d_s-1)}{d(d-1)} - \frac{(d_s-l)^2}{(d-1)^2}\bigg) \\
& = \frac{1-f}{d-1}\bigg( 2fl(1-l)
+ (1-f)\frac{-d_s^2 + d_s((1-2l)(d-1) + 2ld) -2l(1-l)(d-1) - l^2d}{d(d-1)}\bigg) \\
& = \frac{1-f}{d-1}\bigg( 2fl(1-l) 
+ (1-f)\frac{-d_s^2 + d_s(d+2l-1) - dl(2-l)+ 2l(1-l)}{d(d-1)}\bigg) \\
% & < \frac{1-f}{d-1}\bigg( 2fl(1-l) 
% + (1-f)\frac{\frac{1}{4}(d+2l-1)^2 + l^2d}{d(d-1)}\bigg) \\
& \le \frac{1-f}{d-1}\bigg( f\frac{1}{2}
+ (1-f)\frac{\frac{d^2}{4} + \frac{d(2l-1)}{2} - dl(2-l)+ 2l(1-l)}{d(d-1)}\bigg) \\
& = \frac{1-f}{d-1}\bigg( f\frac{1}{2}
+ (1-f)\frac{d^2 - 2d(2l(1-l)+1) + 8l(1-l)}{4d(d-1)}\bigg)
\le \frac{1-f}{d-1}\bigg( f\frac{1}{2}
+ (1-f)\frac{d^2 - 2d}{4d(d-1)}\bigg) \\
& = \frac{1-f}{d-1}\bigg( f\frac{d}{4(d-1)}
+ \frac{d - 2}{4(d-1)}\bigg) 
= \frac{(1-f)d}{4(d-1)^2}\bigg( f + \frac{d - 2}{d}\bigg) \\
& \le \frac{(1-\frac{1}{d})d}{4(d-1)^2}\bigg( \frac{1}{d} + \frac{d - 2}{d}\bigg) 
= \frac{1}{4d}\\
% & < \frac{1}{(d-1)}\bigg( \frac{1}{8} + \frac{2d_s + d_s(d_s-1)}{d} - \frac{(d_s-1)^2}{d-1} \bigg) 
% = \frac{1}{(d-1)}\bigg( \frac{1}{8} + \frac{-d_s^2 + d_s(3d-1) - d}{d(d-1)} \bigg) \\ 
% & < \frac{1}{(d-1)}\bigg( \frac{1}{8} + \frac{-d^2 + d(3d-1) - d}{d(d-1)} \bigg) 
% < \frac{17}{8(d-1)},
\end{align*}

\end{proof}

% Let $f > \frac{1}{d}$ and $l > \frac{d_s}{d}$. Then 
% \begin{align*}
% & \tr(P_0^{\otimes2} \cdot \mbb{E}(\op{u}{u})^{\otimes2}) 
% < l^2f^2 + \frac{2(1-l)^2(d-1)}{d^3} + \frac{4l(1-l)f}{d} + \\
% & + \frac{2(d_s-1)lf}{d} + \frac{2(d_s-1)(1-l)(d-1)}{d^3} 
% + \frac{(d-1)}{d^3}\big( d_s(d_s-1)\big) \\
% % & = 
% % & = l^2\bigg(f^2 + \frac{2(1-f)^2}{d(d-1)} - \frac{4f(1-f)}{d-1} \bigg) 
% % + l\bigg( - \frac{4(1-f)^2}{d(d-1)} + \frac{4f(1-f)}{d-1} 
% %   + \frac{2(d_s-1)f(1-f)}{d-1} - \frac{2(d_s-1)(1-f)^2}{d(d-1)}\bigg) 
% \end{align*}


% \begin{equation}
% \frac{ - \mathtt{d\_s} + 2 l + d \mathtt{d\_s} - 2 d l - \mathtt{d\_s}^{2} + 2 \mathtt{d\_s} f + 2 \mathtt{d\_s} l - 4 f l - 2 l^{2} - 2 d \mathtt{d\_s} f + 2 d f l + l^{2} d + 2 \mathtt{d\_s}^{2} f - f^{2} \mathtt{d\_s} - 4 \mathtt{d\_s} f l + 2 f^{2} l + 4 l^{2} f + 2 d^{2} f l + f^{2} d \mathtt{d\_s} - f^{2} \mathtt{d\_s}^{2} + 2 f^{2} \mathtt{d\_s} l - 2 l^{2} f^{2} - 2 f^{2} d^{2} l - 2 l^{2} d^{2} f - l^{2} f^{2} d + 2 l^{2} f^{2} d^{2}}{\left( -1 + d \right)^{2} d}
% \end{equation}

\bibliographystyle{IEEEtran}
\bibliography{literature.bib}
% \begin{thebibliography}{10}
% \end{thebibliography}


\end{document}
